Considere la ecuación $x^2+y+f(x)y^{\prime}=0$. Puede usar el resultado auxiliar: $xu^{\prime}+u=x$ tiene soluciones $u=x/2+C/x$ para $x\ne0$, y sólo $u=x/2$ se prolonga por el origen. Determine $f$, aplíquelo y estudie la familia ortogonal en los tres apartados siguientes.

\begin{enumerate}[label=(\alph*),ref=\alph*,leftmargin=*]
\item Determine $f$ para que $\mu=x$ sea factor integrante de $x^2+y+f(x)y^{\prime}=0$.
\item Resuelva la ecuación original para $f(x)=1/x+x/2$.
\item Determine las trayectorias ortogonales de la familia de soluciones que corresponde a $f(x)=x/2$.
\end{enumerate}
